\documentclass[../main.tex]{subfiles}

\begin{document}

\chapter{Pendahuluan dan Orientasi OBE}

\begin{subcpmk}
  \item Mahasiswa memahami kerangka orientasi buku ajar Logika Matematika.
  \item Mahasiswa dapat mendeskripsikan tujuan penggunaan paradigma Outcome-Based Education (OBE) dalam proses perkuliahan.
\end{subcpmk}

\section{Tujuan Buku Ajar Logika Matematika}
Buku ajar Logika Matematika ini merupakan pedoman komprehensif yang diperuntukkan bagi mahasiswa S1 Teknik Informatika. Tujuannya adalah untuk meletakkan fondasi \textit{mathematical reasoning} dan penalaran deduktif yang sangat penting dalam:
\begin{itemize}
    \item Membangun kemampuan berpikir kritis (Critical Thinking).
    \item Merancang dan membuktikan validitas suatu algoritma.
    \item Membangun kerangka dasar perangkat keras (arsitektur sistem komputer melalui aljabar Boolean).
    \item Memodelkan masalah dunia nyata ke bahasa pemrograman.
\end{itemize}

\section{Keterkaitan dengan RPS Berbasis OBE}
Di Universitas Bale Bandung, kerangka pencapaian mutu lulusan tertuang di dalam format OBE. Artinya, apa yang dipelajari dan diuji harus memiliki tolok ukur (outcome) yang terdefinisi dengan jelas sejak awal kuliah. Buku ini terhubung langsung dengan dokumen Rencana Pembelajaran Semester (RPS):
\begin{itemize}
    \item Silabus diatur ke dalam 14 unit kompetensi mikro.
    \item Asesmen akhir berpusat pada seberapa mapan mahasiswa merespons Capaian Pembelajaran Lulusan (CPL) dan CPMK, alih-alih sekadar membaca bahan tanpa refleksi.
\end{itemize}

Seluruh bab di buku ini merujuk ke satuan ukur (Sub-CPMK) tersebut untuk memastikan mahasiswa secara eksplisit dapat membuktikan keterampilannya sebelum sesi berakhir.

\begin{aktivitas}
  \item Baca bab "Cara Menggunakan Buku Ini" pada halaman depan dan perhatikan urutan metode evaluasi apa yang akan dijumpai.
\end{aktivitas}

\begin{checklist}
  \item Memahami tujuan dari mata kuliah Logika Matematika.
  \item Mengetahui bahwa setiap materi tertaut ke *Outcome* khusus (Sub-CPMK).
\end{checklist}

\end{document}
